%%%% Originalvorlage von Friedrich Günther, angepasst von Alexander Rogovskyy 

\documentclass[12pt, twocolumn=false, twoside=false, DIV=calc]{scrartcl} %Setup der Dokumentklasse; Schriftgröße, Zweispaltiger Satz (aus), Zweiseitiger Satz (aus), Satzspiegelberechnung (klassischer Buchsatzspiegel)

\usepackage{iftex}
\ifluatex
	\usepackage{fontspec}
\else
	\usepackage[utf8]{inputenc} %Eingabekodierung nach UTF8-Unicode, Erlaubt direkte Eingabe von Umlauten etc.
	\usepackage[T1]{fontenc} %Paket für verbesserten Schriftsupport
\fi


\usepackage[ngerman]{babel} %Neue deutsche Trennreglen, klappt auch mit Lua/Xe

\usepackage{lmodern} %Wechsel zur Schrift "Latin Modern", eine verbesserte Version der Standard-TeX-Schrift "Computer Modern"
\renewcommand{\bfdefault}{b} %Wechsel der Standard-Fettschrift

\usepackage[colorinlistoftodos,prependcaption,textsize=tiny]{todonotes} %TODO's

\usepackage{scrlayer-scrpage} %Zusatzpaket für die Dokumentklasse für Seitenkopf, Seitenfuß, ...
\pagestyle{scrheadings} %Einstellung des Seitenstils
\setkomafont{disposition}{\bfseries} %Überschriften in fett und mit Serifen 

\usepackage{amsmath}
\usepackage{mathtools, amssymb, dsfont, IEEEtrantools} %Pakete für Mathematik-Satz; erweitertes amsmath-Paket, Symbole (griechisches Alphabet etc.), passende doppelt gestrichene Buchstaben, Paket für verbessertes eqnarray
\usepackage{xcolor} %Paket für Verwendung von Farben
\usepackage{graphicx} %Paket zur Einbindung von Grafiken
\usepackage{tikz} %Paket zur Erstellung von Vektorgrafiken und Diagrammen
\usetikzlibrary{cd, babel} %TikZ-Paket "cd" zur Erstellung kommutativer Diagramme; babel ist nötig, damit das Paket in nicht-englischsprachigen Dokumenten funktioniert
\usepackage[colorlinks=true, linktoc=all, linkcolor=red, urlcolor=red]{hyperref} %Paket für die Erstellung von Hyperlinks innerhalb des Dokuments inklusive "Farb-Setup"

\usepackage{enumitem} %Paket für besser konfigurierbare Aufzählungen
\setlist{listparindent=\parindent, parsep=0pt} %Standardeinstellung für Aufzählungen; gleicher Absatzeinzug wie im Dokument, kein zusätzlicher Abstand zwischen einzelnen Punkten

\renewcommand{\labelenumi}{$\mathrm{(\roman{enumi})}$} %Erzwungen aufrechter kleiner römischer Bezeichner für Aufzählungen in erster Hierarchie-Ebene
\renewcommand{\labelenumii}{$\mathrm{(\arabic{enumii})}$} %Erzwungen aufrechter arabischer Bezeichner für Aufzählungen in zweiter Hierarchie-Ebene
\renewcommand{\labelenumiii}{$\mathrm{(\alph{enumiii})}$} %Erzwungen aufrechter kleiner lateinischer Bezeichner für Aufzählungen in dritter Hierarchie-Ebene

\usepackage{nicefrac}
\usepackage[style=german]{csquotes} %Paket für korrekte Anführungszeichen; Verwendung: \enquote{<Satz/Wörter/Wort>}, zum Beispiel \enquote{lustig}.

\clubpenalty10000 %Setzt die Strafe für Schusterjungen und Hurenkindern auf die höchste Stufe, die LaTeX kennt -- verhindert diese Art von typographischen Schweinereien
\widowpenalty10000

\usepackage[amsmath, thmmarks, hyperref]{ntheorem} %Paket zur Definition mathematischer Umgebungen wie Satz, Definition, ...

\theoremstyle{plain}
\theoremheaderfont{\normalfont\bfseries}
\theorembodyfont{\itshape}
\theoremseparator{:}
\theorempreskip{\topsep}
\theorempostskip{\topsep}
\theoremindent0cm
\theoremnumbering{arabic}
%\theoremsymbol{\ensuremath{\square}}
\newtheorem*{Lemma}{Lemma}

\theoremstyle{plain}
\theoremheaderfont{\normalfont\bfseries}
\theorembodyfont{\upshape}
\theoremseparator{:}
\theorempreskip{\topsep}
\theorempostskip{\topsep}
\theoremindent0cm
\theoremnumbering{arabic}
\theoremsymbol{}
\newtheorem{Aufgabe}{Aufgabe}
\newtheorem*{Loesung}{Lösung}
\newtheorem*{Definition}{Definition}

\theoremstyle{plain}
\theoremheaderfont{\normalfont\bfseries}
\theorembodyfont{\upshape}
\theoremseparator{:}
\theorempreskip{\topsep}
\theorempostskip{\topsep}
\theoremindent0cm
\theoremnumbering{arabic}
\theoremsymbol{\ensuremath{\square}}
\newtheorem*{Proof}{Proof}

%---------- Design ---------

% Points in the topright corner
% <dark magic>
\makeatletter
\newcommand{\Repeat}[1]{%
    \expandafter\@Repeat\expandafter{\the\numexpr #1\relax}%
}

\def\@Repeat#1{%
    \ifnum#1>0
        \expandafter\@@Repeat\expandafter{\the\numexpr #1-1\expandafter\relax\expandafter}%
    \else
        \expandafter\@gobble
    \fi
}
\def\@@Repeat#1#2{%
    \@Repeat{#1}{#2}#2%
}
\makeatother
% </dark magic>
\newcounter{headertask}
\setcounter{headertask}{\firsttask{}-1}
\newcommand{\insertpoints}{\begin{tikzpicture}[remember picture,overlay]
	\node[yshift=-7\baselineskip, xshift=0.7\paperwidth] at (current page.north west) [text width=5cm,text centered,above right]{\Large{\begin{tabular}{\Repeat{\numtasks}{c|}c}
		\Repeat{\numtasks}{\stepcounter{headertask}\the\value{headertask} & }
		$\displaystyle\sum$ \\ \hline \Repeat{\numtasks}{&} \end{tabular}}};
\end{tikzpicture}%	
}

% Title, Authors etc.
\newcommand{\inserthead}{
\begin{center}
	\textbf{\course{}} \\
	\emph{\sheet{}}
\end{center}

\noindent\rule{\textwidth}{1pt}
\small{\authors{}}
}


\setcounter{section}{\firsttask{}-1}